\documentclass[11pt]{beamer}

\usetheme{Madrid}
\usecolortheme{default}
\definecolor{unicaucaverde}{RGB}{0,90,60}
\setbeamercolor{structure}{fg=unicaucaverde}
\setbeamertemplate{navigation symbols}{}

\usepackage[utf8]{inputenc}
\usepackage[spanish]{babel}
\usepackage{amsmath,amssymb}

\newcommand{\R}{\mathbb{R}}

\title[Semana 1: Sistemas de ecuaciones]{Sistemas de ecuaciones lineales: introducción y el
  caso de dos incógnitas}
\subtitle{Álgebra Lineal para Ingeniería --- Semana 1 de 16}
\author{Universidad del Cauca}
\date{2026}

\begin{document}

\begin{frame}
  \titlepage
\end{frame}

\begin{frame}{Semana 1: dónde estamos}
  \begin{itemize}
    \item Unidad 1 de 8: Sistemas de ecuaciones lineales.
    \item Subtemas 1.1 y 1.2 --- RA1.1, RA1.2.
    \item 4 horas: teoría + taller de sistemas en general, luego teoría + taller de sistemas
          $2\times 2$.
  \end{itemize}
  \tableofcontents
\end{frame}

\section{Introducción a los sistemas}

\begin{frame}{¿Por qué sistemas de ecuaciones?}
  \begin{itemize}
    \item Leyes de Kirchhoff en circuitos eléctricos.
    \item Equilibrio de fuerzas en armaduras y estructuras.
    \item Balances de materia y energía en procesos.
    \item Precios de equilibrio oferta-demanda.
  \end{itemize}
  \vfill
  En todos los casos: varias incógnitas ligadas por relaciones lineales simultáneas.
\end{frame}

\begin{frame}{Ecuación lineal y sistema}
  \begin{block}{Ecuación lineal}
    \[
      a_1x_1+a_2x_2+\cdots+a_nx_n = b
    \]
  \end{block}
  \begin{block}{Sistema de $m$ ecuaciones con $n$ incógnitas}
    \[
    \begin{aligned}
      a_{11}x_1+\cdots+a_{1n}x_n &= b_1 \\
      &\;\vdots \\
      a_{m1}x_1+\cdots+a_{mn}x_n &= b_m
    \end{aligned}
    \]
  \end{block}
  Una \textbf{solución} es una $n$-tupla que satisface las $m$ ecuaciones a la vez.
\end{frame}

\begin{frame}{Consistente / inconsistente}
  \begin{itemize}
    \item \textbf{Consistente}: tiene al menos una solución.
    \item \textbf{Inconsistente}: no tiene ninguna.
  \end{itemize}
  \begin{alertblock}{Un sistema consistente tiene...}
    exactamente una solución, o infinitas --- nunca un número finito mayor que uno. Se ve ya en
    el caso $2\times 2$ y se demuestra en general en la Semana 2.
  \end{alertblock}
  \vfill
  \textbf{Taller (hora 2):} verificar soluciones; traducir un problema de circuito o mezcla a un
  sistema.
\end{frame}

\section{Dos ecuaciones con dos incógnitas}

\begin{frame}{Sistema $2\times 2$: dos rectas}
  \[
  \begin{aligned}
    a_1x+b_1y &= c_1 \\
    a_2x+b_2y &= c_2
  \end{aligned}
  \]
  Cada ecuación es una recta en $\R^2$. Resolver el sistema = hallar la intersección.
\end{frame}

\begin{frame}{Los tres casos}
  \begin{columns}
    \begin{column}{0.33\textwidth}
      \centering \textbf{Secantes}\\ solución única
    \end{column}
    \begin{column}{0.33\textwidth}
      \centering \textbf{Paralelas}\\ sin solución
    \end{column}
    \begin{column}{0.33\textwidth}
      \centering \textbf{Coincidentes}\\ infinitas soluciones
    \end{column}
  \end{columns}
  \vfill
  \begin{exampleblock}{Solución única}
    $x+y=5,\; x-y=1 \;\Rightarrow\; (3,2)$
  \end{exampleblock}
\end{frame}

\begin{frame}{Sustitución y eliminación}
  \begin{block}{Sustitución}
    Despejar una variable de una ecuación y sustituir en la otra.
  \end{block}
  \begin{block}{Eliminación}
    Combinar (sumar/restar múltiplos de) las ecuaciones para cancelar una variable. Generaliza
    directamente a Gauss-Jordan (Semana 2).
  \end{block}
\end{frame}

\begin{frame}{Ejemplos: los tres casos}
  \textbf{Sin solución:}
  \[
    2x+y=4, \qquad 4x+2y=3
  \]
  (mismas pendientes, distinto término independiente: rectas paralelas)

  \vfill
  \textbf{Infinitas soluciones:}
  \[
    x-2y=3, \qquad -2x+4y=-6
  \]
  (la segunda ecuación es $-2$ veces la primera: la misma recta)

  \vfill
  \textbf{Taller (hora 4):} \texttt{visualizacion/} sección \emph{1.1--1.2} --- variar
  $a_1,b_1,c_1,a_2,b_2,c_2$ y observar los tres casos.
\end{frame}

\section{Soluciones de ejercicios}

\begin{frame}{Ej. 1 — Forma general $3 \times 4$}
  Sistema de $m=3$ ecuaciones con $n=4$ incógnitas ($x_1,x_2,x_3,x_4$):
  \[
  \begin{aligned}
    a_{11}x_1 + a_{12}x_2 + a_{13}x_3 + a_{14}x_4 &= b_1 \\
    a_{21}x_1 + a_{22}x_2 + a_{23}x_3 + a_{24}x_4 &= b_2 \\
    a_{31}x_1 + a_{32}x_2 + a_{33}x_3 + a_{34}x_4 &= b_3
  \end{aligned}
  \]
  \begin{block}{Elementos}
    $a_{ij}$ son los coeficientes, $b_i$ los términos independientes.
  \end{block}
\end{frame}

\begin{frame}{Ej. 2 — Verificación de solución}
  ¿Es $(2,-1,1)$ solución de $\{x_1+x_2+x_3=2,\; 2x_1-x_2+x_3=6\}$?
  \vfill
  \begin{itemize}
    \item Ecuación 1: $2 + (-1) + 1 = 2$ \quad \textcolor{unicaucaverde}{\checkmark (Cumple)}
    \item Ecuación 2: $2(2) - (-1) + 1 = 4 + 1 + 1 = 6$ \quad \textcolor{unicaucaverde}{\checkmark (Cumple)}
  \end{itemize}
  \vfill
  \begin{alertblock}{Conclusión}
    Satisface ambas ecuaciones $\Rightarrow$ \textbf{Sí es solución}.
  \end{alertblock}
\end{frame}

\begin{frame}{Ej. 3 — Aplicación en circuitos}
  Modelado: $3I_1 - I_2 = 5$ y $-2I_1 + 5I_2 = 0$.
  \vfill
  \textbf{Resolviendo por eliminación:}
  \begin{align*}
    5 \times \text{Eq}_1 + \text{Eq}_2 \implies 13I_1 = 25 \implies I_1 = \frac{25}{13} \text{ A}
  \end{align*}
  Sustituyendo $I_1$: $5I_2 = 2I_1 \implies I_2 = \frac{10}{13} \text{ A}$.
  \vfill
  \begin{block}{Resultado}
    \[ \boxed{I_1 = \frac{25}{13}\text{ A}, \quad I_2 = \frac{10}{13}\text{ A}} \]
  \end{block}
\end{frame}

\begin{frame}{Ej. 4 — Sustitución y clasificación}
  Sistema: $\{3x-y=7,\; x+2y=0\}$.
  \vfill
  De la 2da ecuación: $x = -2y$. Sustituyendo en la 1ra:
  \[
    3(-2y) - y = 7 \implies -7y = 7 \implies y = -1 \implies x = 2
  \]
  \vfill
  \begin{block}{Resultado}
    \[ \boxed{(x,y) = (2,-1)} \]
    \textbf{Sistema consistente determinado} (solución única, rectas secantes).
  \end{block}
\end{frame}

\begin{frame}{Ej. 5 — Análisis de coeficientes}
  Sistema: $\{x+3y=6,\; 2x+6y=10\}$.
  \vfill
  \begin{itemize}
    \item Pendientes: $m_1 = -1/3$, $m_2 = -2/6 = -1/3$ (paralelas).
    \item Comparación de razones: $\frac{1}{2} = \frac{3}{6} \neq \frac{6}{10}$.
  \end{itemize}
  \vfill
  \begin{alertblock}{Resultado}
    Rectas paralelas distintas $\Rightarrow$ \textbf{Sin solución} (inconsistente).
  \end{alertblock}
\end{frame}

\begin{frame}{Ej. 6 — Construcción de sistema indeterminado}
  Primera ecuación dada: $3x + y = 9$.
  \vfill
  \begin{block}{Construcción}
    Multiplicando la ecuación por un escalar no nulo ($c = 2$):
    \[
      2(3x + y) = 2(9) \implies 6x + 2y = 18
    \]
  \end{block}
  \vfill
  El sistema $\{3x+y=9,\; 6x+2y=18\}$ representa la misma recta y tiene \textbf{infinitas soluciones}.
\end{frame}

\begin{frame}{Cierre de la semana}
  \begin{itemize}
    \item RA1.1: forma general de un sistema, verificación de soluciones.
    \item RA1.2: sistemas $2\times 2$, sustitución/eliminación, tres casos geométricos.
    \item Control formativo al inicio de la Semana 2.
  \end{itemize}
  \vfill
  \centering \Large Semana 2: $m$ ecuaciones con $n$ incógnitas --- Gauss-Jordan y sistemas
  homogéneos.
\end{frame}

\end{document}
